\subsection{Restriction and Corestriction}

Let G and \(G'\) be groups, F be an abelian group, and \(\tau\) be a homomorphism from G to the automorphism group of F. Let \(u: G' \to G\) be a group homomorphism. We write \(\tau' = \tau \circ u\) .

If \(\psi : \mathrm{G} \times \mathrm{G} \to \mathrm{F}\) is a 2-cocycle of \(\mathrm{G}\) with values in \(\mathrm{F}\) , then the mapping \(\psi \circ (u \times u)\) from \(\mathrm{G}' \times \mathrm{G}'\) to \(\mathrm{F}\) given by \((g_1, g_2) \mapsto \psi(u(g_1), u(g_2))\) is a 2-cocycle of \(\mathrm{G}'\) with values in \(\mathrm{F}\) , and the mapping \(\psi \mapsto \psi \circ (u \times u)\) from \(\mathrm{Z}^2(\mathrm{G}, \mathrm{F})\) to \(\mathrm{Z}^2(\mathrm{G}', \mathrm{F})\) induces a group homomorphism \(u^*: \mathrm{H}^2(\mathrm{G}, \mathrm{F}) \to \mathrm{H}^2(\mathrm{G}', \mathrm{F})\) . For any \(\lambda \in \mathrm{H}^2(\mathrm{G}, \mathrm{F})\) , the element \(u^*(\lambda)\) is called the inverse image of \(\lambda\) by \(u\) . When \(\mathrm{G}'\) is a subgroup of \(\mathrm{G}\) and \(u: \mathrm{G}' \to \mathrm{G}\) is the canonical injection, the homomorphism \(u^*\) is called the restriction homomorphism from \(\mathrm{G}\) to \(\mathrm{G}'\) ; we denote it by \(\operatorname{Res}_{\mathrm{G}'}^{\mathrm{G}}\) . When \(\mathrm{G}\) is a quotient of \(\mathrm{G}'\) and \(u: \mathrm{G}' \to \mathrm{G}\) is the canonical surjection, the homomorphism \(u^*\) is called the inflation homomorphism from \(\mathrm{G}\) to \(\mathrm{G}'\) .

PROPOSITION 7. --- In the notation from above, the following diagram commutes:

\[
\begin{array}{c} \operatorname{Ex} _ {\tau} (\mathrm{G}, \mathrm{F}) \xrightarrow {u ^ {*}} \operatorname{Ex} _ {\tau} (\mathrm{G} ^ {\prime}, \mathrm{F}) \\ \Biggl \downarrow \Theta_ {\tau} \qquad \qquad \qquad \Biggl \downarrow \Theta_ {\tau^ {\prime}} \\ \operatorname{H} ^ {2} (\mathrm{G}, \mathrm{F}) \xrightarrow {u ^ {*}} \operatorname{H} ^ {2} (\mathrm{G} ^ {\prime}, \mathrm{F})  . \end{array}
\]

Let H be a subgroup of G of finite index. Let s be a section of the canonical surjection from G to \(H\backslash G\). We denote by \((g,x)\mapsto x\cdot g\) the right action of G on \(H\backslash G\) induced by the right action of G on itself by multiplication. Note that for every \(x\in H\backslash G\) and \(g\in G\) , the element \(s(x)gs(x\cdot g)^{-1}\) belongs to H. For any mapping \(c: \mathrm{H} \times \mathrm{H} \to \mathrm{F}\) , we can therefore define a mapping \(\widetilde{c}_s: \mathrm{G} \times \mathrm{G} \to \mathrm{F}\) by the relation

\[
\widetilde {c} _ {s} (g _ {1}, g _ {2}) = \prod_ {x \in \mathrm{H} \backslash \mathrm{G}} ^ {s (x) ^ {- 1}} c \bigl (s (x) g _ {1} s (x \cdot g _ {1}) ^ {- 1}, s (x \cdot g _ {1}) g _ {2} s (x \cdot g _ {1} g _ {2}) ^ {- 1} \bigr)
\]

for \(g_{1}, g_{2} \in G\) . The mapping \(c \mapsto \widetilde{c}_{s}\) is a group homomorphism from \(F^{H \times H}\) to \(F^{G \times G}\) .

Lemma 4. --- If \(c: \mathrm{H} \times \mathrm{H} \to \mathrm{F}\) is a 2-cocycle of H with values in F, then \(\widetilde{c}_s\) is a 2-cocycle of G with values in F.

Let \(g_1, g_2\) , and \(g_3\) be elements of G. For every \(i \in \{1, 2, 3\}\) , we define a mapping \(h_i: \mathrm{H} \backslash \mathrm{G} \to \mathrm{H}\) by the relation

\[
h _ {i} (x) = s \left(x \cdot g _ {1} \dots g _ {i - 1}\right) g _ {i} s \left(x \cdot g _ {1} \dots g _ {i}\right) ^ {- 1}
\]

for \(x\in \mathrm{H}\backslash \mathrm{G}\) . Note that

\[
\begin{array}{l} {h _ {1} (x) h _ {2} (x) = s (x) g _ {1} g _ {2} s (x \cdot g _ {1} g _ {2}) ^ {- 1}} \\ {h _ {2} (x) h _ {3} (x) = s (x \cdot g _ {1}) g _ {2} g _ {3} s (x \cdot g _ {1} g _ {2} g _ {3}) ^ {- 1}} \end{array}
\]

for \(x \in \mathrm{H} \backslash \mathrm{G}\) . We then have the relations

\[
\begin{array}{r l} & ^ {g _ {1}} \widetilde {c} _ {s} (g _ {2}, g _ {3}) \widetilde {c} _ {s} (g _ {1} g _ {2}, g _ {3}) ^ {- 1} \widetilde {c} _ {s} (g _ {1}, g _ {2} g _ {3}) \widetilde {c} _ {s} (g _ {1}, g _ {2}) ^ {- 1} \\ & \quad = \prod_ {x \in \mathrm{H} \backslash \mathrm{G}} ^ {g _ {1} s (x) ^ {- 1}} c \big (s (x) g _ {2} s (x \cdot g _ {2}) ^ {- 1}, s (x \cdot g _ {2}) g _ {3} s (x \cdot g _ {2} g _ {3}) ^ {- 1} \big) \\ & \qquad \times \prod_ {x \in \mathrm{H} \backslash \mathrm{G}} ^ {s (x) ^ {- 1}} \Big (c (h _ {1} (x) h _ {2} (x), h _ {3} (x)) ^ {- 1} \\ & \qquad \qquad \qquad \qquad \times c (h _ {1} (x), h _ {2} (x) h _ {3} (x))   c (h _ {1} (x), h _ {2} (x)) ^ {- 1} \Big) \\ & = \prod_ {x \in \mathrm{H} \backslash \mathrm{G}} ^ {s (x \cdot g _ {1} ^ {- 1}) ^ {- 1} h _ {1} (x \cdot g _ {1} ^ {- 1})} c (h _ {2} (x \cdot g _ {1} ^ {- 1}), h _ {3} (x \cdot g _ {1} ^ {- 1})) \\ & \qquad \times \prod_ {x \in \mathrm{H} \backslash \mathrm{G}} ^ {s (x) ^ {- 1} h _ {1} (x)} c (h _ {2} (x), h _ {3} (x)) ^ {- 1} \\ & = 1, \end{array}
\]

where the second equality follows from the fact that c is a 2-cocycle.

Lemma 5. --- If \(c\) is a 2-coboundary, then so is \(\widetilde{c}_s\) .

Let \(t: \mathrm{H} \to \mathrm{F}\) be a mapping such that \(c = \partial t\) . Let \(\widetilde{t}_s: \mathrm{G} \to \mathrm{F}\) be the mapping defined by

\[
\widetilde {t} _ {s} (g) = \prod_ {x \in \mathrm{H} \backslash \mathrm{G}} ^ {s (x) ^ {- 1}} t (s (x) g s (x \cdot g) ^ {- 1})
\]

for \(g \in G\) . It suffices to prove that \(\widetilde{c}_{s} = \partial\widetilde{t}_{s}\) , which follows from the relations

\[
\begin{array}{l} \widetilde {c} _ {s} (g _ {1}, g _ {2}) = \prod_ {x \in \mathrm{H} \backslash \mathrm{G}} ^ {s (x) ^ {- 1} s (x) g _ {1} s (x \cdot g _ {1}) ^ {- 1}} t (s (x \cdot g _ {1}) g _ {2} s (x \cdot g _ {1} g _ {2}) ^ {- 1}) \\ \qquad \times \prod_ {x \in \mathrm{H} \backslash \mathrm{G}} ^ {s (x) ^ {- 1}} \big (t (s (x) g _ {1} g _ {2} s (x \cdot g _ {1} g _ {2}) ^ {- 1}) ^ {- 1} t (s (x) g _ {1} s (x \cdot g _ {1}) ^ {- 1}) \big) \\ = \partial \widetilde {t} _ {s} (g _ {1}, g _ {2}) \end{array}
\]

for \(g_{1}, g_{2} \in G\) .

Lemma 6. --- Let c be a 2-cocycle of H with values in F. The image of \(\widetilde{c}_{s}\) in the group \(\mathrm{H}^{2}(\mathrm{G},\mathrm{F})\) does not depend on the choice of the section s.

Let \(s'\) be a section of the canonical surjection from G to \(H\backslash G\). Let \(h: \mathrm{H} \backslash \mathrm{G} \to \mathrm{H}\) be the mapping characterized by the relation

\[
s ^ {\prime} (x) = h (x) s (x)
\]

for \(x \in H \backslash G\) . The 2-cocycle \(\widetilde{c}_{s'}\) is then given by the relations

\[
\begin{array}{l} \widetilde {c} _ {s ^ {\prime}} (g _ {1}, g _ {2}) \\ = \prod_ {x \in \mathrm{H} \backslash \mathrm{G}} ^ {s (x) ^ {- 1} h (x) ^ {- 1}} c \big (h (x) s (x) g _ {1} s ^ {\prime} (x \cdot g _ {1}) ^ {- 1}, s ^ {\prime} (x \cdot g _ {1}) g _ {2} s ^ {\prime} (x \cdot g _ {1} g _ {2}) ^ {- 1} \big) \\ = \prod_ {x \in \mathrm{H} \backslash \mathrm{G}} ^ {s (x) ^ {- 1}} c \big (s (x) g _ {1} s (x \cdot g _ {1}) ^ {- 1} h (x \cdot g _ {1}) ^ {- 1}, h (x \cdot g _ {1}) s (x \cdot g _ {1}) g _ {2} s ^ {\prime} (x \cdot g _ {1} g _ {2}) ^ {- 1} \big) \\ \qquad \times \prod_ {x \in \mathrm{H} \backslash \mathrm{G}} ^ {s (x) ^ {- 1}} \Big (c \big (h (x) ^ {- 1}, s ^ {\prime} (x) g _ {1} g _ {2} s ^ {\prime} (x \cdot g _ {1} g _ {2}) ^ {- 1} \big) ^ {- 1} \\ \qquad \qquad \qquad \qquad \qquad \qquad \qquad \qquad \qquad \qquad \qquad \qquad \qquad \qquad \qquad \qquad \qquad \qquad \qquad \qquad \qquad \qquad \qquad \qquad \qquad \qquad \qquad \qquad \qquad \qquad \qquad \qquad \qquad \qquad \\ = \prod_ {x \in \mathrm{H} \backslash \mathrm{G}} ^ {g _ {1} s (x \cdot g _ {1}) ^ {- 1}} c \big (h (x \cdot g _ {1}) ^ {- 1}, h (x \cdot g _ {1}) s (x \cdot g _ {1}) g _ {2} s ^ {\prime} (x \cdot g _ {1} g _ {2}) ^ {- 1} \big) \\ \qquad \times \prod_ {x \in \mathrm{H} \backslash \mathrm{G}} ^ {s (x) ^ - 1} c \big (s (x) g _ {1} s (x \cdot g _ {1}) ^ {- 1}, s (x \cdot g _ {1}) g _ {2} s ^ {\prime} (x \cdot g _ {1} g _ {2}) ^ {- 1} \big) \\ \qquad \times \prod_ {x \in \mathrm{H} \backslash \mathrm{G}} ^ {s (x) ^ {- 1}} c \big (s (x) g _ {1} s (x \cdot g _ {1}) ^ {- 1}, h (x \cdot g _ {1}) ^ {- 1} \big) ^ {- 1} \\ \qquad \times \prod_ {x \in \mathrm{H} \backslash \mathrm{G}} ^ {s (x) ^ {- 1}} \Big (c \big (h (x) ^ {- 1}, s ^ {\prime} (x) g _ {1} g _ {2} s ^ {\prime} (x \cdot g _ {1} g _ {2}) ^ {- 1} \big) ^ {{- 1}} \\ \qquad \qquad \qquad \qquad \qquad \qquad \qquad \qquad \qquad \qquad \qquad \qquad \times c (h (x) ^ {- 1}, s ^ {\prime} (x) g _ {1} s ^ {\prime} (x \cdot g _ {1}) ^ {- 1})   {\Big)} \end{array}
\]

\[
\begin{array}{l} = \prod_ {x \in \mathrm{H} \backslash \mathrm{G}} ^ {s (x) ^ {- 1}} c \big (s (x) g _ {1} s (x \cdot g _ {1}) ^ {- 1}, s (x \cdot g _ {1}) g _ {2} s ^ {\prime} (x \cdot g _ {1} g _ {2}) ^ {- 1} \big) \\ \qquad \times \prod_ {x \in \mathrm{H} \backslash \mathrm{G}} ^ {s (x) ^ {- 1}} c \big (s (x) g _ {1} s (x \cdot g _ {1}) ^ {- 1}, h (x \cdot g _ {1}) ^ {- 1} \big) ^ {- 1} \\ \qquad \times \prod_ {x \in \mathrm{H} \backslash \mathrm{G}} ^ {g _ {1} s (x) ^ {- 1}} c \big (h (x) ^ {- 1}, h (x) s (x) g _ {2} s ^ {\prime} (x \cdot g _ {2}) ^ {- 1} \big) \\ \qquad \times \prod_ {x \in \mathrm{H} \backslash \mathrm{G}} ^ {s (x) ^ {- 1}} \Big (c \big (h (x) ^ {- 1}, s ^ {\prime} (x) g _ {1} g _ {2} s ^ {\prime} (x \cdot g _ {1} g _ {2}) ^ {- 1} \big) ^ {- 1} \\ \qquad \qquad \qquad \qquad \qquad \qquad \qquad \qquad \qquad \qquad \qquad \qquad \qquad \qquad \qquad \qquad \qquad \qquad \times c \big (h (x) ^ {- 1}, s ^ {\prime} (x) g _ {1} s ^ {\prime} (x \cdot g _ {1}) ^ {- 1} \big) \Big) \end{array}
\]

for all \(g_1, g_2 \in \mathbf{G}\) . The first equality comes from the cocycle relation (VIII, p.~295, formula (7)) applied to the elements

\[
h (x) ^ {- 1}, \quad h (x) s (x) g _ {1} s ^ {\prime} (x \cdot g _ {1}) ^ {- 1}, \quad \mathrm{and} \quad s ^ {\prime} (x \cdot g _ {1}) g _ {2} s ^ {\prime} (x \cdot g _ {1} g _ {2}) ^ {- 1};
\]

the second is obtained by applying the cocycle relation to the elements

\[
s (x) g _ {1} s \left(x \cdot g _ {1}\right) ^ {- 1}, \quad h \left(x \cdot g _ {1}\right) ^ {- 1}, \quad \text { and } \quad h \left(x \cdot g _ {1}\right) s \left(x \cdot g _ {1}\right) g _ {2} s ^ {\prime} \left(x \cdot g _ {1} g _ {2}\right) ^ {- 1};
\]

and the last simply uses the fact that the mapping \(x \mapsto x \cdot g_{1}\) is a permutation of \(H\backslash G\).

The last two lines of the obtained expression correspond to a 2-coboundary. We find that \(\widetilde{c}_{s'}\) has the same class in \(\mathrm{H}^2 (\mathrm{G},\mathrm{F})\) as the cocycle whose value in \((g_{1},g_{2})\in \mathrm{G}^{2}\) is given by the expression

\[
\begin{array}{l} \prod_ {x \in \mathrm{H} \backslash \mathrm{G}} ^ {s (x) ^ {- 1}} c \big (s (x) g _ {1} s (x \cdot g _ {1}) ^ {- 1}, s (x \cdot g _ {1}) g _ {2} s (x \cdot g _ {1} g _ {2}) ^ {- 1} h (x \cdot g _ {1} g _ {2}) ^ {- 1} \big) \\ \qquad \qquad \qquad \times \prod_ {x \in \mathrm{H} \backslash \mathrm{G}} ^ {s (x) ^ {- 1}} c \big (s (x) g _ {1} s (x \cdot g _ {1}) ^ {- 1}, h (x \cdot g _ {1}) ^ {- 1} \big) ^ {- 1} \\ = \prod_ {x \in \mathrm{H} \backslash \mathrm{G}} ^ {g _ {1} s (x \cdot g _ {1}) ^ {- 1}} c \big (s (x \cdot g _ {1}) g _ {2} s (x \cdot g _ {1} g _ {2}) ^ {- 1}, h (x \cdot g _ {1} g _ {2}) ^ {- 1} \big) ^ {- 1} \\ \qquad \qquad \times \prod_ {x \in \mathrm{H} \backslash \mathrm{G}} ^ {s (x) ^ {- 1}} c \big (s (x) g _ {1} s (x \cdot g _ {1}) ^ {- 1}, s (x \cdot g _ {1}) g _ {2} s (x \cdot g _ {1} g _ {2}) ^ {- 1} \big) \\ \qquad \qquad \times \prod_ {x \in \mathrm{H} \backslash \mathrm{G}} ^ {s (x) ^ {- 1}} \Big (c \big (s (x) g _ {1} g _ {2} s (x \cdot g _ {1} g _ {2}) ^ {- 1}, h (x \cdot g _ {1} g _ {2}) ^ {- 1} \big) \\ \qquad \qquad \qquad \times c \big (s (x) g _ {1} s (x \cdot g _ {1}) ^ {- 1}, h (x \cdot g _ {1}) ^ {- 1} \big) ^ {- 1} \Big) \end{array}
\]

\[
\begin{array}{l} = \prod_ {x \in \mathrm{H} \backslash \mathrm{G}} ^ {s (x) ^ {- 1}} c \big (s (x) g _ {1} s (x \cdot g _ {1}) ^ {- 1}, s (x \cdot g _ {1}) g _ {2} s (x \cdot g _ {1} g _ {2}) ^ {- 1} \big) \\ \quad \times \prod_ {x \in \mathrm{H} \backslash \mathrm{G}} ^ {g _ {1} s (x) ^ {- 1}} c \big (s (x) g _ {2} s (x \cdot g _ {2}) ^ {- 1}, h (x \cdot g _ {2}) ^ {- 1} \big) ^ {- 1} \\ \quad \times \prod_ {x \in \mathrm{H} \backslash \mathrm{G}} ^ {s (x) ^ {- 1}} \Big (c \big (s (x) g _ {1} g _ {2} s (x \cdot g _ {1} g _ {2}) ^ {- 1}, h (x \cdot g _ {1} g _ {2}) ^ {- 1} \big) \\ \qquad \qquad \qquad \qquad \qquad \qquad \qquad \qquad \qquad \qquad \qquad \qquad \qquad \times c \big (s (x) g _ {1} s (x \cdot g _ {1}) ^ {- 1}, h (x \cdot g _ {1}) ^ {- 1} \big) ^ {- 1} \Big), \end{array}
\]

where the first equality follows from the cocycle relation applied to the elements

\[
s (x) g _ {1} s (x \cdot g _ {1}) ^ {- 1}, \quad s (x \cdot g _ {1}) g _ {2} s (x \cdot g _ {1} g _ {2}) ^ {- 1}, \quad \text { and } \quad h (x \cdot g _ {1} g _ {2}) ^ {- 1}.
\]

The last two lines of the obtained expression correspond to a 2-coboundary. We find that the class of \(\widetilde{c}_{s'}\) coincides with that of \(\widetilde{c}_s\) .

We have constructed a homomorphism from the group \(\mathrm{H}^{2}(\mathrm{H},\mathrm{F})\) to the group \(\mathrm{H}^{2}(\mathrm{G},\mathrm{F})\) , which we call the corestriction homomorphism from H to G and denote by \(Cor_{H}^{G}\) .

PROPOSITION 8. --- Let H be a subgroup of G of finite index. The endomorphism \(Cor_{H}^{G} \circ Res_{H}^{G}\) of the group \(H^{2}(G, F)\) coincides with multiplication by the index \((G : H)\) .

Let \(\alpha\) be an element of \(\mathrm{H}^2 (\mathrm{G},\mathrm{F})\) , and let \(c\) be an element of \(\mathrm{Z}^2 (\mathrm{G},\mathrm{F})\) representing \(\alpha\) . The element \(\mathrm{Cor}_{\mathrm{H}}^{\mathrm{G}}\circ \mathrm{Res}_{\mathrm{H}}^{\mathrm{G}}(\alpha)\) is the class of the cocycle whose value in \((g_{1},g_{2})\in \mathrm{G}^{2}\) is given by the expression

\[
\begin{array}{l} \prod_ {x \in \mathrm{H} \backslash \mathrm{G}} ^ {s (x) ^ {- 1}} c (s (x) g _ {1} s (x \cdot g _ {1}) ^ {- 1}, s (x \cdot g _ {1}) g _ {2} s (x \cdot g _ {1} g _ {2}) ^ {- 1}) \\ = \prod_ {x \in \mathrm{H} \backslash \mathrm{G}} c (g _ {1} s (x \cdot g _ {1}) ^ {- 1}, s (x \cdot g _ {1}) g _ {2} s (x \cdot g _ {1} g _ {2}) ^ {- 1}) \\ \quad \times \prod_ {x \in \mathrm{H} \backslash \mathrm{G}} \big (c (s (x) ^ {- 1}, s (x) g _ {1} g _ {2} s (x \cdot g _ {1} g _ {2}) ^ {- 1}) ^ {- 1} c (s (x) ^ {- 1}, s (x) g _ {1} s (x \cdot g _ {1}) ^ {- 1}) \big) \\ = \prod_ {x \in \mathrm{H} \backslash \mathrm{G}} ^ {g _ {1}} c (s (x \cdot g _ {1}) ^ {- 1}, s (x \cdot g _ {1}) g _ {2} s (x \cdot g _ {1} g _ {2}) ^ {- 1}) \\ \quad \times \prod_ {x \in \mathrm{H} \backslash \mathrm{G}} c (g _ {1}, g _ {2} s (x \cdot g _ {1} g _ {2}) ^ {- 1}) c (g _ {1}, s (x \cdot g _ {1}) ^ {- 1}) ^ {- 1} \\ \quad \times \prod_ {x \in \mathrm{H} \backslash \mathrm{G}} \big (c (s (x) ^ {- 1}, s (x) g _ {1} g _ {2} s (x \cdot g _ {1} g _ {2}) ^ {- 1}) ^ {- 1} c (s (x) ^ {- 1}, s (x) g _ {1} s (x \cdot g _ {1}) ^ {- 1}) \big) \end{array}
\]

\[
\begin{array}{l} = \prod_ {x \in \mathrm{H} \backslash \mathrm{G}} c (g _ {1}, g _ {2} s (x \cdot g _ {1} g _ {2}) ^ {- 1}) c (g _ {1}, s (x \cdot g _ {1}) ^ {- 1}) ^ {- 1} \\ \times \prod_ {x \in \mathrm{H} \backslash \mathrm{G}} ^ {g _ {1}} c (s (x) ^ {- 1}, s (x) g _ {2} s (x g _ {2}) ^ {- 1}) \\ \times \prod_ {x \in \mathrm{H} \backslash \mathrm{G}} \left(c (s (x) ^ {- 1}, s (x) g _ {1} g _ {2} s (x \cdot g _ {1} g _ {2}) ^ {- 1}) ^ {- 1} c (s (x) ^ {- 1}, s (x) g _ {1} s (x \cdot g _ {1}) ^ {- 1})\right). \end{array}
\]

The first equality comes from the cocycle relation (VIII, p.~295, formula (7)) applied to the elements

\[
s (x) ^ {- 1}, \quad s (x) g _ {1} s (x \cdot g _ {1}) ^ {- 1}, \quad \mathrm{and} \quad s (x \cdot g _ {1}) g _ {2} s (x \cdot g _ {1} g _ {2}) ^ {- 1};
\]

the second is obtained by applying the cocycle relation to the elements

\[
g _ {1}, \quad s (x \cdot g _ {1}) ^ {- 1}, \quad \mathrm{and} \quad s (x \cdot g _ {1}) g _ {2} s (x \cdot g _ {1} g _ {2}) ^ {- 1}.
\]

By removing a coboundary, we obtain that \(\mathrm{Cor}_{\mathrm{H}}^{\mathrm{G}} \circ \mathrm{Res}_{\mathrm{H}}^{\mathrm{G}}(\alpha)\) is the class of the cocycle whose value in \((g_{1}, g_{2}) \in \mathrm{G}^{2}\) is given by the expression

\[
\begin{array}{l} \prod_ {x \in \mathrm{H} \backslash \mathrm{G}} c (g _ {1}, g _ {2} s (x \cdot g _ {1} g _ {2}) ^ {- 1}) c (g _ {1}, s (x \cdot g _ {1}) ^ {- 1}) ^ {- 1} \\ = \prod_ {x \in \mathrm{H} \backslash \mathrm{G}} ^ {g _ {1}} c (g _ {2}, s (x \cdot g _ {1} g _ {2}) ^ {- 1}) ^ {- 1} c (g _ {1} g _ {2}, s (x \cdot g _ {1} g _ {2}) ^ {- 1}) c (g _ {1}, g _ {2}) \\ \times \prod_ {x \in \mathrm{H} \backslash \mathrm{G}} c (g _ {1}, s (x \cdot g _ {1}) ^ {- 1}) ^ {- 1} \\ = \prod_ {x \in \mathrm{H} \backslash \mathrm{G}} ^ {g _ {1}} c (g _ {2}, s (x \cdot g _ {2}) ^ {- 1}) ^ {- 1} c (g _ {1} g _ {2}, s (x \cdot g _ {1} g _ {2}) ^ {- 1}) c (g _ {1}, s (x \cdot g _ {1}) ^ {- 1}) ^ {- 1} \\ \times c (g _ {1}, g _ {2}) ^ {\text {(G:H)}}, \end{array}
\]

which proves the result.

